\section{Methods}
\label{sec:methods}

\subsection{Features}

All features at issue time $t$ use only information available at $t$: station reports at or
before $t$, and fire detections acquired at or before $t-L$, where $L$ is the fire-data latency.

\paragraph{Station state (C).}
Current visibility, minimum visibility over the last 6\,h and 24\,h, current haze code, share of
reports with a haze code over 24\,h and 72\,h, share with precipitation over 24\,h, relative
humidity, temperature, wind speed, \textbf{wind direction} (sine and cosine), local hour, day of
year, and the station-month climatological event rate (estimated from training years only).
Including the wind direction in the baseline means that a station's typical ``haze wind'' is
already available to every model. A directional fire feature can only help by locating the fires
relative to that wind.

\paragraph{Fire counts (D).}
Detections within 100\,km in the last 24\,h, and within 300\,km in the last 72\,h.

\paragraph{Fire exposure.}
For detections $i$ within 300\,km acquired in the 72\,h before $t-L$, let $d_i$ be the
great-circle distance to the station, $a_i$ the detection age, $\mathrm{FRP}_i$ its fire radiative
power (1\,MW when missing), $\hat w$ the unit vector of the wind's flow direction (the reported
meteorological ``from'' direction plus 180$^\circ$), and $\hat h_i$ the unit vector from the
detection towards the station. We compare a non-directional exposure
\begin{equation}
  S^{\mathrm{iso}} = \sum_i \mathrm{FRP}_i\, e^{-d_i/\lambda}\, e^{-a_i/\tau}
\end{equation}
with a directional (upwind) exposure
\begin{equation}
  S^{\mathrm{dir}} = \sum_i \mathrm{FRP}_i\, e^{-d_i/\lambda}\, \max(0, \hat w\cdot\hat h_i)\, e^{-a_i/\tau},
\end{equation}
with $\lambda=50$\,km and $\tau=24$\,h as in the deployed Pavana model (not tuned here). The only
difference between the two is the alignment term.

\paragraph{Placebos.}
A placebo should keep everything about $S^{\mathrm{dir}}$ except its link to the actual wind--fire
geometry. Rotating the wind by 180$^\circ$ or 90$^\circ$ is a natural candidate, but it is not a
clean negative control. A rotated exposure is still computed from the current wind and the actual
fire positions, so it carries geometric information correlated with $S^{\mathrm{dir}}$ (for
example $S^{\mathrm{dir}}+S^{180}=\sum_i w_i\,|\hat w\cdot\hat h_i|$, a measure of how aligned the
fires are with the wind axis). It also shares the missingness pattern of $S^{\mathrm{dir}}$
exactly. We therefore report rotations only for contrast, and use two families of placebos that
break the current wind--fire geometry:
\begin{itemize}
  \item \emph{Time-shifted wind} ($S^{\mathrm{shift}}_k$): the same station's wind from $k$ hours
        earlier, with $k\in\{72, 168, 336, 720\}$; $k=168$\,h is the pre-specified main placebo
        (same hour of day). A shifted wind can still share the regime of the current wind
        (monsoon flow persists for weeks) and has its own missingness pattern, so we report how
        similar the shifted and current winds are (Appendix~\ref{app:extra}).
  \item \emph{Mask-preserving placebos}, which have exactly the missingness of $S^{\mathrm{dir}}$:
        directions permuted at random among the hours with a defined wind of the same station
        and calendar month ($S^{\mathrm{perm}}$, three seeds), and the 168\,h-shifted direction
        (the last defined value at or before $t-168$\,h) applied only where the current
        direction is defined ($S^{\mathrm{shift,m}}_{168}$).
\end{itemize}

\paragraph{Wind-report categories.}
$S^{\mathrm{dir}}$ is undefined whenever the selected report has no usable direction. This
happens for three different reasons with different meanings: calm wind (\texttt{00000KT},
\pv{audit.ALL.calm}\% of issue hours), variable wind (\texttt{VRB}, \pv{audit.ALL.vrb}\%), and a
missing wind group (\pv{audit.ALL.wmiss}\%). To separate missingness from direction we use two
non-directional controls. $S^{\mathrm{iso}}_{\mathrm{calm}}$ is the non-directional exposure set
to missing whenever the direction is undefined (the exact missingness of $S^{\mathrm{dir}}$). The
category control (``cat'') gives the model indicators for the three categories and the
non-directional exposure split by category ($S^{\mathrm{iso}}\cdot\mathbb 1[\text{category}]$),
so any interaction between exposure and reporting state is available explicitly.

\paragraph{Persistent flow (exploratory).}
Smoke arriving over 72\,h follows the persistent flow rather than one hourly surface report. We
therefore also compute $S^{\mathrm{dir}}$ with the 24\,h speed-weighted vector-mean wind (defined
in 95\% of hours), with its own shifted placebo.

\subsection{Variants}

\begin{table}[h]
\centering
\small
\begin{tabular}{ll}
\toprule
Variant & Features \\
\midrule
C & station state \\
D & C + fire counts \\
D0 & D + $S^{\mathrm{iso}}$ \\
\textbf{E} & D0 + $S^{\mathrm{dir}}$ \\
\textbf{Rshift} & D0 + $S^{\mathrm{shift}}_{168}$ (pre-specified main placebo) \\
Rshift$k$ & D0 + $S^{\mathrm{shift}}_k$, $k\in\{72,336,720\}$\,h \\
Rperm1--3 & D0 + $S^{\mathrm{perm}}$ (mask-preserving, three seeds) \\
Rshiftm168 & D0 + $S^{\mathrm{shift,m}}_{168}$ (mask-preserving) \\
R180, R90 & D0 + $S^{\mathrm{dir}}$ with rotated wind (not clean nulls; for contrast) \\
D0m & D0 + $S^{\mathrm{iso}}_{\mathrm{calm}}$ (missingness of $S^{\mathrm{dir}}$, no direction) \\
D0cat, Ecat & D0 / E + wind-report categories and category-split exposure \\
D0v, Ev, Rshiftv, D0mv & as D0, E, Rshift, D0m with vegetation fires only (FIRMS type 0) \\
Ep, Rpshift & D0 + persistent-flow $S^{\mathrm{dir}}$ / its shifted placebo (exploratory) \\
E$_\text{served}$ & D + $S^{\mathrm{dir}}$ (the deployed feature set) \\
\bottomrule
\end{tabular}
\caption{LightGBM feature sets. The directional question is E against the placebos and the
non-directional controls (D0m, D0cat), not E against D0 alone. Logistic-regression versions of C,
D0, E, Rshift, Ep, Rpshift, D0cat, Ecat and Rperm1 (prefix ``lg'', with $\log(1+x)$ exposures and
missingness indicators) repeat the key comparisons in a second model family. Variants that were
not in the frozen protocol (all except C, D, D0, E, R180, R90 and E$_\text{served}$) are listed
with their dates in Appendix~\ref{app:protocol}.}
\label{tab:variants}
\end{table}

Every variant is a LightGBM binary classifier \citep{lightgbm} with the deployed hyperparameters
(learning rate 0.05, 31 leaves, at least 200 rows per leaf, feature fraction 0.9, bagging fraction
0.8, $\ell_2=1$) and a fixed budget of 400 boosting rounds, so no variant receives more tuning than
another. The random seed controls bagging and feature subsampling, and single fits turned out to
differ by almost 0.1 AP between seeds in some test years (Section~\ref{sec:seeds}). Each LightGBM
variant is therefore a five-seed ensemble: the mean predicted probability of five fits that
differ only in the seed. The protocol specified a single fit; those results are reported
separately. We also report three reference models: persistence (the event indicator at $t$),
station-month climatology, and logistic regression on C.

\subsection{Evaluation}

\paragraph{Leave-one-year-out (primary).}
Each calendar year 2015--2025 is held out in turn. The model is trained on all other years,
excluding 14 days on either side of the test year to limit leakage across boundaries. Pooled
out-of-fold predictions are scored once per variant and lead time. Training on later years to
predict an earlier one is not an operational forecast. It is a way to score several independent
haze seasons, which is what the attribution question needs.

\paragraph{Operational split (secondary).}
Train on 2015--2020 and test on 2023--2025, which is the configuration of the deployed model.

\paragraph{Latency.}
The main setting is $L=3$\,h, the nominal latency of FIRMS near-real-time data. Sensitivity runs
use $L\in\{0,6,12,24\}$\,h, for D0, E and Rshift at 24\,h (five-seed ensembles).

\paragraph{Metrics.}
Events are rare (0.2--3.7\% of hours per station), so the primary metric is average precision
(AP, area under the precision--recall curve). We also report ROC-AUC, the Brier skill score
against station-month climatology, and the share of event hours captured when alerting on the top
2\% of hours.

\paragraph{Pooled and per-season scores.}
Pooled out-of-fold AP is dominated by the largest season (2015), whose fold is also the hardest
extrapolation. We therefore also report the mean of within-year $\Delta$AP over the years with at
least 100 event hours (ten of eleven years), which weights each season equally.

\paragraph{Uncertainty.}
Differences in AP are assessed with a paired block bootstrap (1000 resamples) over
station~$\times$~ISO-week blocks. This keeps within-week autocorrelation and the pairing of
variants on identical rows. We report 95\% percentile intervals. A haze episode often affects
several airports at once, so station-level blocks can understate uncertainty. As a sensitivity
check we repeat the key contrasts with blocks that hold all stations together: ISO weeks,
calendar months, and calendar years (500 resamples each). Eleven years do not make eleven
independent episodes, and the year-block intervals, from only eleven clusters, are coarse.

\paragraph{Hypotheses.}
H1 (primary): $\Delta\mathrm{AP}(\mathrm{E}-\mathrm{D0})>0$ at $h=24$\,h, $L=3$\,h, under
leave-one-year-out. H2: the directional gain is larger at 24\,h than at 6\,h. H3: placebos do not
reproduce the gain. H4: the directional gain shrinks as $L$ grows. All results are reported
whatever their sign. Appendix~\ref{app:protocol} lists the deviations from the frozen protocol, in
order: the shifted-wind placebo and missingness control, the persistent-flow variant, the
logistic family and per-season scores, the controls added in response to review (wind-report
categories, mask-preserving placebos, further shifts, vegetation-only fires, regional bootstrap
blocks), and the five-seed ensembles. All were added after first results and are labelled as
such.
